\section{Importance of Discourse in Live Commentary}
\label{sec:importance_of_discourse}

The role of shoutcasters during a live esports event is not limited to providing real-time descriptions of what is happening in the game. Viewers are inherently limited in the amount of information they are able to obtain from the observers, due to the fact that full control of the in-game camera is only available to producers and team staff \cite{cheung2011starcraft}. Therefore, shoutcasters are responsible for providing the audience with further information that they would not be able to obtain otherwise.

\medskip 

In addition to that, shoutcasters have to take into account the difference in the level of expertise from members of the audience. Some events might have to be explained and interpreted for newer viewers, while also providing more in-depth analysis for more experienced ones. In both of these cases, it's necessary for them to consider what, and how to supply explanations in the context of the match at hand \cite{penney2021shoutcasters}. 

\medskip

Moreover, similar to traditional sports broadcasting, building a narrative around the match through the use of stories from the game's history has shown to be greatly appreciated by audiences \cite{lee2014automated,smith1997storytellers}. These discourses are ideal to compensate for moments of low action during games, helping to build emotional engagement with the audience, and foster emotional investment to the players and teams \cite{ryan1993narrative}. 
\medskip

As a summary of the previously stated points, we have identified four main discursive functions of shoutcasting for LoL esports events:

\begin{itemize}
  \item \textbf{Play-by-Play}: refers to real-time narration of ongoing action in the game, usually during high-action moments which have a high influence in the outcome of the match.
  \item \textbf{Strategic Analysis}: refers to commentary made by shoutcasters with more knowledge about the game, providing insights and explanations towards events of the ongoing match, usually during lower-action moments, or as a follow-up to high-action moments.
  \item \textbf{Narrative Storytelling}: refers to the use of stories, and theories, made by the casters, to build a narrative around the match in order to increase the expectations around it.
  \item \textbf{Emotional and \textit{Hype}}: refers to the use of language and tone as the result of the casters' own emotional engagement with the match, which looks to be conveyed to the audience, commonly as a follow-up to high-action moments, or at the start and end of a match.
\end{itemize}

\medskip

Through these functions, the complex gameplay becomes more manageable for audience members, and slow moments of the game become opportunities for engagement, either through analysis or entertaining storytelling. 

\medskip

While this categorization is not exhaustive, it provides a structured foundation for the computational modeling of our research. Both identifying and collecting examples to understand each of them is necessary for systematic analysis, and the possibility of developing tools for shoutcaster assistance in the future. 